\documentclass[11pt]{article}

\usepackage{amsmath}
\usepackage{amsthm}
\usepackage{amsfonts}
\usepackage{hyperref}	
\usepackage{xcolor}	
\usepackage{graphicx}
\usepackage{multirow}

\setlength{\textwidth}     {16.0cm}
\setlength{\textheight}    {21.0cm}
\setlength{\evensidemargin}{ 0.0cm}
\setlength{\oddsidemargin} { 0.0cm}
\setlength{\topmargin}     {-0.5cm}
\setlength{\baselineskip}  { 0.7cm}

\newcommand{\R}{\mathbb{R}}
\newcommand{\N}{\mathbb{N}}
\newcommand{\xkjtrial}{x^{k,j}_{\mathrm{trial}}}
\newcommand{\minimize}{\mathop{\text{Minimize}}}
\newcommand{\flow}{f_{\mathrm{low}}}
\newcommand{\simplex}{\textstyle\sum}

\newtheorem{theorem}{Theorem}[section]
\newtheorem{corollary}{Corollary}[section]
\newtheorem{lemma}{Lemma}[section]
\newtheorem{definition}{Definition}[section]
\newtheorem{assumption}{Assumption}
\renewcommand\theassumption{A\arabic{assumption}}

% Show line numbers
\usepackage[mathlines]{lineno}
\renewcommand\linenumberfont{\normalfont\small\color{gray}}
\linenumbers

% Show labels
\usepackage{showlabels}


\begin{document}

\title{A Quasi-Newton-like regularized approach\\ to the order-value optimization problem\thanks{This work has been funded by the Vice-Rectorate for Research, Outreach and Social Projection of the University of the Atlantic through the 2022 Internal Call for Proposals to Strengthen Research Groups with Tools for Creation, Innovation, and Research.}}

\author{
    Carlos Araujo\thanks{Mathematics Program, Faculty of Basic Sciences, University of Atlantic, Puerto Colombia, Atlántico, Colombia. e-mail: carlosaraujo@uniatlantico.edu.co}
    \and
    Miguel Caro\thanks{Mathematics Program, Faculty of Basic Sciences, University of Atlantic, Puerto Colombia, Atlántico, Colombia. e-mail: miguelcaro@uniatlantico.edu.co}
    \and
    Gustavo Quintero\thanks{Mathematics Program, Faculty of Basic Sciences, University of Atlantic, Puerto Colombia, Atlántico, Colombia. e-mail: gdquintero@uniatlantico.edu.co}
}

% \date{June 18, 2025\footnote{Revision made on \today.}}
\date{\today}
\maketitle

\begin{abstract}
The order-value function of order $p$ returns, at each point, the $p$th
smallest value among a finite family of functions $f_1,\dots,f_m$;
minimizing it recovers the minimization of the minimum and of the maximum
of the family in the extreme cases $p=1$ and $p=m$, and provides a flexible
tool for robust estimation, in which a prescribed number of outlying
observations are discarded. We propose a regularized method for minimizing
the order-value function subject to box constraints. At each iteration the
method minimizes a strongly convex model that combines a piecewise-linear
term built from the active gradients, a quasi-Newton curvature term defined
by a symmetric positive semidefinite matrix, and a quadratic regularization;
the curvature term carries second-order information without requiring explicit
second derivatives, and a recently proposed first-order regularized method is
recovered as the special case in which this term vanishes. We introduce an
approximate optimality condition for the problem and prove that the method is
well defined and attains the same worst-case iteration- and
evaluation-complexity bounds as its first-order counterpart, so that the
incorporation of curvature preserves the theoretical guarantees while offering
the potential for higher-quality steps in practice. The applicability of the
method to parameter-estimation problems with outliers is illustrated.\\

\noindent
\textbf{Keywords:} Order-value optimization, regularized models,
quasi-Newton methods, complexity, algorithms, applications.\\

\noindent
\textbf{Mathematics Subject Classification:} 90C30, 65K05, 49M37, 90C60, 68Q25.\\
\end{abstract}

\section{Introduction} \label{sec:intro}

Many optimization problems involve objective functions that are obtained not
by aggregating a collection of simpler functions, but by selecting one of them
according to its rank. The prototypical example is the order-value function of
order $p$, introduced in \cite{dunder1,dunder2}: given $m$ functions
$f_1,\dots,f_m$, its value at a point $x$ is the $p$th smallest among
$f_1(x),\dots,f_m(x)$. Functions of this kind belong to the broader class of
generalized order-value functions---those whose value at $x$ is governed by
order relations among the elements of a set $\{f_i(x)\}_{i\in I}$---a class
systematized in \cite{mtop}. Since the index that attains the $p$th position
may switch abruptly as $x$ moves, the order-value function is generally
nonsmooth even when each $f_i$ is continuously differentiable.

The extreme choices $p=1$ and $p=m$ recover the minimum and the maximum of
$\{f_i(x)\}$, respectively, so the order-value function interpolates between
these two classical aggregates; the intermediate cases are the ones of
practical interest. When $x$ encodes portfolio positions and $f_i(x)$ is the
loss incurred under scenario $i$, the order-value function coincides with the
discrete value-at-risk (VaR) measure, a standard risk measure in financial
risk management \cite{jorion}. When $f_i(x)$ measures the misfit of a
parametric model to the $i$th observation, minimizing the $p$th order-value
function fits the parameters while discarding the $m-p$ worst-fitted
observations. This makes order-value optimization a natural tool for robust
estimation, where one seeks to remain insensitive to a minority of grossly
corrupted data, or outliers.

The order-value optimization (OVO) problem was first formulated in
\cite{dunder1}, where a steepest-descent-type method was proposed and shown to
converge to points satisfying a weak optimality condition. Sharper optimality
conditions and a reformulation as a nonlinear program with equilibrium
constraints were subsequently developed in \cite{dunder2}. A quasi-Newton
method generalizing the scheme of \cite{dunder1} was introduced in
\cite{yano1}, and a global optimization strategy combining multistart with a
tunneling technique was proposed in \cite{salva}. On the
computational-complexity side, \cite{jiang} established that minimizing the
order-value function over a polytope is NP-hard in the strong sense.

The systematic study of worst-case complexity in continuous optimization
originates in the cubic-regularization analysis of \cite{nesterov}, and
complexity guarantees have since been obtained for a broad range of problems
(see, e.g., \cite{bgmst2,bgmst1} and the monograph \cite{cogtbook}). Within
this line, \cite{alvarezovo} recently proposed a first-order method for
box-constrained OVO that minimizes, at each iteration, a quadratically
regularized piecewise-linear model of $f$, and established iteration- and
evaluation-complexity bounds for reaching an approximate optimality condition.

In the present work we generalize the approach of \cite{alvarezovo} by
enriching the regularized model with second-order information. At each
iteration our method minimizes a model built from the same piecewise-linear
convex term, a curvature term governed by a symmetric positive semidefinite
matrix $B$, and a Tikhonov-type quadratic regularization. The matrix $B$ is
updated in a quasi-Newton fashion and therefore requires no explicit second
derivatives, while the regularization renders the model strongly convex, so
that each subproblem has a unique minimizer. Setting $B=0$ recovers the
first-order method of \cite{alvarezovo} exactly, so the proposed method is a
strict generalization of it. We derive an approximate optimality condition for
the box-constrained problem and prove that the method is well defined and
attains the \emph{same} worst-case iteration- and evaluation-complexity bounds
as the first-order method. Thus the incorporation of curvature does not degrade
the theoretical guarantees, while it offers the potential for higher-quality
steps, and hence faster convergence, in practice---a behavior we illustrate on
parameter-estimation problems in the presence of outliers.

The remainder of the paper is organized as follows. Section~\ref{sec:method}
defines the box-constrained order-value optimization problem and the proposed
regularized method. Section~\ref{sec:complexity} introduces an appropriate
optimality condition and establishes the well-definiteness, convergence, and
complexity of the method. Section~\ref{sec:numerical} reports illustrative
numerical experiments on parameter estimation in the presence of outliers.
Final remarks are given in Section~\ref{sec:final}.

\medskip
\noindent\textit{Notation.} The symbol $\|\cdot\|$ denotes the Euclidean norm.
For $i=1,\dots,n$, $e_i\in\mathbb{R}^{n}$ denotes the $i$th column of the

\section{A quasi-Newton regularized method} \label{sec:method}

Let $f_i : \mathbb{R}^n \to \mathbb{R}$, $i = 1, \ldots, m$, be given and let
$p \in \{1, \ldots, m\}$ be fixed. The \textit{order-value function of order
$p$} is the mapping $f : \mathbb{R}^n \to \mathbb{R}$ defined by
\begin{equation}
    f(x) = f_{i_p(x)}(x), \label{ovo-function}
\end{equation}
where $\{i_1(x), \ldots, i_m(x)\}$ is a permutation of $\{1, \ldots, m\}$ such
that
\begin{equation}
    f_{i_1(x)}(x) \leq f_{i_2(x)}(x) \leq \cdots \leq f_{i_m(x)}(x);
    \label{fi-sort}
\end{equation}
that is, $f(x)$ is the $p$th smallest value among $f_1(x), \ldots, f_m(x)$.
The boundary cases $p = 1$ and $p = m$ recover, respectively, the minimum and
the maximum of $\{f_1(x), \ldots, f_m(x)\}$.

If $f_1, \ldots, f_m$ are continuous, then so is $f$. Differentiability,
however, is not inherited: $f$ may be nonsmooth at a point $x$ at which the
$p$th smallest value is attained by more than one index, i.e.\ whenever there
exist $i \neq j$ with $f_i(x) = f_j(x) = f(x)$, since the active index $i_p(x)$
may then change as $x$ varies.

We address the box-constrained order-value optimization problem
\begin{equation}
    \minimize_{x \in \Omega} f(x), \label{ovo-problem}
\end{equation}
where $\Omega := \{x \in \mathbb{R}^n \mid \ell \leq x \leq u\}$ with
$\ell, u \in \mathbb{R}^n$ and $\ell_i < u_i$ for $i = 1, \ldots, n$. Box
constraints arise naturally in parameter estimation, where the unknowns are
confined to physically or statistically meaningful bounds.

For $\delta > 0$ and $x \in \Omega$, define the \textit{active index set}
\begin{equation}
    I_\delta(x) := \{i \in \{1, \ldots, m\} \mid
    f(x) - \delta \leq f_i(x) \leq f(x) + \delta\}, \label{idelta}
\end{equation}
the set of indices whose values lie within a band of half-width $\delta$ around
$f(x)$. As $\delta \to 0$, $I_\delta(x)$ shrinks to the set of indices that
attain the order-value,
\begin{equation*}
    I_0(x) := \{i \in \{1, \ldots, m\} \mid f_i(x) = f(x)\}.
\end{equation*}

The method builds, around a reference point $\bar{x} \in \Omega$, a local model
of $f$ that combines the first-order information of the active components with
curvature carried by a symmetric positive semidefinite matrix
$B \in \mathbb{R}^{n \times n}$. In a quasi-Newton framework, $B$ is updated
from successive gradients, so that no second derivatives are computed
explicitly. For given $\bar{x} \in \Omega$, $B \succeq 0$, $\delta > 0$, and
$\sigma > 0$, the regularized model is
\begin{equation}
    \Psi(x; \bar{x}, B, \delta, \sigma) =
    \max_{i \in I_\delta(\bar{x})}
    \left\lbrace \nabla f_i(\bar{x})^T(x - \bar{x}) 
    + \frac{1}{2} \left[ (x - \bar{x})^T B (x - \bar{x})
    + \sigma \|x - \bar{x}\|^2 \right]\right\rbrace. \label{model}
\end{equation}
The first term is the piecewise-linear convex function that captures the
worst-case first-order behavior of $f$ over the active components; the bracket
adds the quasi-Newton curvature $(x - \bar{x})^T B (x - \bar{x})$ and a
Tikhonov regularization $\sigma \|x - \bar{x}\|^2$. Since $\sigma > 0$, the
model is strongly convex even when $B$ is only positive semidefinite, so the
subproblem \eqref{ovo-subpro} has a unique minimizer. For $B = 0$, the model
reduces to the first-order regularized model of \cite{alvarezovo}, which the
present method thus generalizes.

\medskip
\noindent\textbf{Algorithm~\ref{sec:method}.1.} Let $\delta > 0$,
$\sigma_{\min} > 0$, $M > 0$, $\alpha \in (0,1)$, $\gamma > 1$, and
$x^0 \in \Omega$ be given. Initialize $k \leftarrow 0$.
\begin{description}
\item[Step~1.] Initialize $j \leftarrow 0$ and $\sigma_{k,j} = \sigma_{\min}$.
\item[Step~2.] Choose a symmetric positive semidefinite matrix
$B_{k,j} \in \mathbb{R}^{n \times n}$ with $\|B_{k,j}\| \leq M$ and compute
$\xkjtrial$ as the global minimizer of
    \begin{equation} \label{ovo-subpro}
    \minimize_{x \in \Omega}\, \Psi(x; x^k, B_{k,j}, \delta, \sigma_{k,j}).
    \end{equation}
\item[Step~3.] If
    \begin{equation} \label{ovo-armijo}
    f(\xkjtrial) \leq f(x^k) - \alpha \|\xkjtrial - x^k\|^2
    \end{equation}
does not hold, set $\sigma_{k,j+1} = \gamma \sigma_{k,j}$,
$j \leftarrow j + 1$, and go to Step~2.
\item[Step~4.] Set $x^{k+1} = \xkjtrial$, $\sigma_k = \sigma_{k,j}$,
$j_k = j$, $k \leftarrow k+1$, and go to Step~1.
\end{description}

The bound $\|B_{k,j}\| \leq M$ keeps the curvature from dominating the model.
The freedom in choosing $B_{k,j}$ lets the method accommodate standard
quasi-Newton updates such as BFGS or SR1, and the choice $B_{k,j} = 0$ for all
$k, j$ recovers the first-order method of \cite{alvarezovo}.

The requirements on $\xkjtrial$ are mild. The theory only needs $\xkjtrial$ to
be a stationary point of \eqref{ovo-subpro} with
$\Psi(\xkjtrial; x^k, B_{k,j}, \delta, \sigma_{k,j}) \leq 0$; since
$\Psi(x^k; x^k, B_{k,j}, \delta, \sigma_{k,j}) = 0$, any descent method started
at $x^k$ produces such iterates. Moreover, \eqref{ovo-subpro} minimizes a
strongly convex function over the compact convex set $\Omega$, so it has a
unique global minimizer and every stationary point is global; standard convex
solvers apply directly.

\section{Optimality condition and complexity} \label{sec:complexity}

\begin{assumption} \label{a1}
The functions $f_1, \ldots, f_m$ belong to $\mathcal{C}^1(\Omega)$; that is,
they are continuously differentiable on $\Omega$.
\end{assumption}

We first show that a local minimizer of \eqref{ovo-problem} also minimizes the
model $\Psi$ built at the corresponding reference point. These results connect
the original nonsmooth problem with the convex subproblems solved at each
iteration and underlie the optimality condition introduced below.

\begin{theorem} \label{teo:conv1}
Suppose that Assumption~\ref{a1} holds. If $x^*$ is a local minimizer of
\eqref{ovo-problem}, then it is also a local minimizer of
$\minimize_{x \in \Omega}\, \Psi(x; x^*, 0, 0, 0)$.
\end{theorem}

\begin{proof}
Suppose not. Then there is $y^* \in \Omega$ with
$\Psi(y^*; x^*, 0, 0, 0) < \Psi(x^*; x^*, 0, 0, 0) = 0$, i.e.\
$\nabla f_i(x^*)^T(y^* - x^*) < 0$ for all $i \in I_0(x^*)$. By
Assumption~\ref{a1}, for each such $i$,
\[
\lim_{t \to 0^+} \frac{f_i(x^* + t(y^* - x^*)) - f_i(x^*)}{t}
= \nabla f_i(x^*)^T(y^* - x^*) < 0,
\]
so there is $\bar{t}_i > 0$ with $f_i(x^* + t(y^* - x^*)) < f_i(x^*) = f(x^*)$
for $t \in (0, \bar{t}_i]$. With $\bar{t} = \min_{i \in I_0(x^*)} \bar{t}_i$,
\begin{equation} \label{bajaste}
f_i(x^* + t(y^* - x^*)) < f(x^*) \quad
\text{for all } i \in I_0(x^*),\ t \in (0, \bar{t}].
\end{equation}
For $j \notin I_0(x^*)$, continuity of $f_1, \ldots, f_m$ gives, for $t$ small,
\begin{align}
f_i(x^* + t(y^* - x^*)) &< f_j(x^* + t(y^* - x^*))
&&\text{if } i \in I_0(x^*),\ f(x^*) < f_j(x^*), \label{arriba}\\
f_i(x^* + t(y^* - x^*)) &> f_j(x^* + t(y^* - x^*))
&&\text{if } i \in I_0(x^*),\ f(x^*) > f_j(x^*). \label{abajo}
\end{align}
Hence the ranking positions of the active indices $I_0(x^*)$ within
$\{f_i(x^* + t(y^* - x^*))\}_{i=1}^m$ are preserved for $t$ small, so the $p$th
position is still occupied by some $i \in I_0(x^*)$ and
$f(x^* + t(y^* - x^*)) = f_i(x^* + t(y^* - x^*))$. With \eqref{bajaste},
$f(x^* + t(y^* - x^*)) < f(x^*)$ for $t$ small. Since $\Omega$ is convex,
$x^* + t(y^* - x^*) \in \Omega$, contradicting the local minimality of $x^*$.
\end{proof}

\begin{theorem} \label{teo:conv2}
Suppose that Assumption~\ref{a1} holds and let $B \succeq 0$ and $\sigma \geq 0$.
If $x^*$ is a local minimizer of \eqref{ovo-problem}, then it is also a local
minimizer of $\minimize_{x \in \Omega}\, \Psi(x; x^*, B, 0, \sigma)$.
\end{theorem}

\begin{proof}
Suppose not: there is $y^* \in \Omega$ with $\Psi(y^*; x^*, B, 0, \sigma) < 0$,
i.e.
\[
\max_{i \in I_0(x^*)} \{ \nabla f_i(x^*)^T(y^* - x^*) \}
+ \tfrac{1}{2} \big[ (y^* - x^*)^T B (y^* - x^*) + \sigma \|y^* - x^*\|^2 \big]
< 0.
\]
Since $B \succeq 0$ and $\sigma \geq 0$, the bracket is nonnegative, so
$\max_{i \in I_0(x^*)} \{ \nabla f_i(x^*)^T(y^* - x^*) \} < 0$, that is,
$\Psi(y^*; x^*, 0, 0, 0) < 0$, contradicting Theorem~\ref{teo:conv1}.
\end{proof}

\begin{theorem} \label{teo:conv3}
Suppose that Assumption~\ref{a1} holds and let $B \succeq 0$ and
$\sigma, \delta \geq 0$. If $x^*$ is a local minimizer of \eqref{ovo-problem},
then it is also a local minimizer of
$\minimize_{x \in \Omega}\, \Psi(x; x^*, B, \delta, \sigma)$.
\end{theorem}

\begin{proof}
Suppose not: there is $y^* \in \Omega$ with
$\Psi(y^*; x^*, B, \delta, \sigma) < 0$. Since
$I_0(x^*) \subseteq I_\delta(x^*)$, we have
$\max_{i \in I_0(x^*)}\{\cdot\} \leq \max_{i \in I_\delta(x^*)}\{\cdot\}$, hence
\[
\max_{i \in I_0(x^*)} \{ \nabla f_i(x^*)^T(y^* - x^*) \}
+ \tfrac{1}{2} \big[ (y^* - x^*)^T B (y^* - x^*) + \sigma \|y^* - x^*\|^2 \big]
\leq \Psi(y^*; x^*, B, \delta, \sigma) < 0,
\]
that is, $\Psi(y^*; x^*, B, 0, \sigma) < 0$, contradicting
Theorem~\ref{teo:conv2}.
\end{proof}

Throughout, $\Sigma_q := \{ \lambda \in \mathbb{R}^q \mid
\sum_{j=1}^q \lambda_j = 1,\ \lambda_j \geq 0 \}$ denotes the unit simplex.
Following \cite{alvarezovo}, the next consequence of Theorem~\ref{teo:conv3}
provides a first-order necessary condition for \eqref{ovo-problem} and
motivates the stopping criterion below.

\begin{corollary} \label{cor:exact}
Suppose that Assumption~\ref{a1} holds and let $\delta \geq 0$. If $x^*$ is a
local minimizer of \eqref{ovo-problem}, then there exist
$\mu \in \Sigma_{|I_\delta(x^*)|}$ and $\nu^\ell, \nu^u \in \mathbb{R}^n_+$ such
that $(\ell_i - x^*_i)\nu^\ell_i = (x^*_i - u_i)\nu^u_i = 0$ for
$i = 1, \ldots, n$ and
\begin{equation} \label{exact-opt}
\sum_{i \in I_\delta(x^*)} \mu_i \nabla f_i(x^*)
+ \sum_{i=1}^n (\nu^u_i - \nu^\ell_i) e_i = 0.
\end{equation}
\end{corollary}

\begin{proof}
By Theorem~\ref{teo:conv3} with $B = 0$ and $\sigma = 0$, $x^*$ minimizes
$x \mapsto \max_{i \in I_\delta(x^*)} \{ \nabla f_i(x^*)^T(x - x^*) \}$ over
$\Omega$. Applying the KKT conditions for the minimization of a maximum of
differentiable functions over a box \cite[Lemma~3.1]{alvarezovo} yields
multipliers $\mu \in \Sigma_{|I_\delta(x^*)|}$ and
$\nu^\ell, \nu^u \in \mathbb{R}^n_+$ with the stated complementarity and
\eqref{exact-opt}.
\end{proof}

\begin{definition} \label{def:optcond}
Given $\epsilon > 0$ and $\delta > 0$, we say that $x \in \Omega$ satisfies the
\textit{approximate optimality condition} $C_{\delta,\epsilon}$ if there exist
$\mu \in \Sigma_{|I_\delta(x)|}$ and $\nu^\ell, \nu^u \in \mathbb{R}^n_+$ such
that $(\ell_i - x_i)\nu^\ell_i = (x_i - u_i)\nu^u_i = 0$ for $i = 1, \ldots, n$
and
\begin{equation} \label{miaprox-box}
\left\| \sum_{i \in I_\delta(x)} \mu_i \nabla f_i(x)
+ \sum_{i=1}^{n} (\nu^u_i - \nu^\ell_i) e_i \right\| \leq \epsilon.
\end{equation}
\end{definition}

Condition $C_{\delta,\epsilon}$ requires the gradient combination of the active
components, corrected by the box multipliers, to be small in norm. For
$\epsilon = 0$ it reduces to the exact condition \eqref{exact-opt} of
Corollary~\ref{cor:exact}. It serves as the stopping criterion for
Algorithm~\ref{sec:method}.1.

\begin{corollary} \label{cor:step}
Suppose that Assumption~\ref{a1} holds. Then, for every $k$ and
$j = 0, \ldots, j_k$, there exist $\mu \in \Sigma_{|I_\delta(x^k)|}$ and
$\nu^\ell, \nu^u \in \mathbb{R}^n_+$ such that
$(\ell_i - [\xkjtrial]_i)\nu^\ell_i = ([\xkjtrial]_i - u_i)\nu^u_i = 0$ for
$i = 1, \ldots, n$ and
\begin{equation} \label{eqcoro1}
\xkjtrial - x^k = \big( B_{k,j} + \sigma_{k,j} I \big)^{-1}
\left[ \sum_{i=1}^{n} (\nu^\ell_i - \nu^u_i) e_i
- \sum_{i \in I_\delta(x^k)} \mu_i \nabla f_i(x^k) \right].
\end{equation}
\end{corollary}

\begin{proof}
Subproblem \eqref{ovo-subpro} minimizes the strongly convex function
$\Psi(\cdot\,; x^k, B_{k,j}, \delta, \sigma_{k,j})$ over the box $\Omega$. By
the KKT conditions for a maximum of differentiable functions over a box
\cite[Lemma~3.1]{alvarezovo}, there exist $\mu \in \Sigma_{|I_\delta(x^k)|}$ and
$\nu^\ell, \nu^u \in \mathbb{R}^n_+$ with the stated complementarity such that
\[
\sum_{i \in I_\delta(x^k)} \mu_i
\big[ \nabla f_i(x^k) + (B_{k,j} + \sigma_{k,j} I)(\xkjtrial - x^k) \big]
+ \sum_{i=1}^n (\nu^u_i - \nu^\ell_i) e_i = 0,
\]
where $\nabla f_i(x^k) + (B_{k,j} + \sigma_{k,j} I)(\xkjtrial - x^k)$ is the
gradient at $\xkjtrial$ of the $i$th component of $\Psi$. Since
$\sum_i \mu_i = 1$, the curvature term factors out; as $\sigma_{k,j} > 0$ and
$B_{k,j} \succeq 0$, the matrix $B_{k,j} + \sigma_{k,j} I$ is positive definite,
and solving for $\xkjtrial - x^k$ gives \eqref{eqcoro1}.
\end{proof}

\begin{assumption} \label{a2}
For all $k$ and $j = 0, \ldots, j_k$, the multipliers $\nu^\ell, \nu^u$ of
Corollary~\ref{cor:step} satisfy
$\max_{i = 1, \ldots, n} |\nu^\ell_i - \nu^u_i| \leq c_\nu$ for some constant
$c_\nu > 0$ independent of $k$ and $j$.
\end{assumption}

The Mangasarian--Fromovitz constraint qualification guarantees that, for each
$(k,j)$, these multipliers are bounded \cite{gauvin}; Assumption~\ref{a2}
strengthens this to a single bound uniform in $k$ and $j$, which is natural
since $\Omega$ is compact and the active gradients are bounded under
Assumption~\ref{a1}.

\begin{assumption} \label{a3}
There is a constant $c_\nabla > 0$ such that $\|\nabla f_i(x)\| \leq c_\nabla$
for all $i = 1, \ldots, m$ and all $x \in \Omega$.
\end{assumption}

Assumption~\ref{a3} holds automatically under Assumption~\ref{a1}, since
$\Omega$ is bounded. Combining Corollary~\ref{cor:step} with
Assumptions~\ref{a2} and~\ref{a3} bounds the step. Since $B_{k,j} \succeq 0$,
the eigenvalues of $B_{k,j} + \sigma_{k,j} I$ are at least $\sigma_{k,j}$, so
$\|(B_{k,j} + \sigma_{k,j} I)^{-1}\| \leq 1/\sigma_{k,j}$, and \eqref{eqcoro1}
gives
\[
\|\xkjtrial - x^k\| \leq \frac{1}{\sigma_{k,j}}
\left( \sum_{i=1}^n |\nu^\ell_i - \nu^u_i|
+ \sum_{i \in I_\delta(x^k)} \mu_i \|\nabla f_i(x^k)\| \right)
\leq \frac{n\,c_\nu + c_\nabla}{\sigma_{k,j}},
\]
where the last inequality uses Assumptions~\ref{a2}, \ref{a3} and
$\mu \in \Sigma_{|I_\delta(x^k)|}$. Thus
\begin{equation} \label{ccotasub}
\|\xkjtrial - x^k\| \leq \frac{c_x}{\sigma_{k,j}}, \qquad
c_x := n\,c_\nu + c_\nabla.
\end{equation}
The step shrinks as $\sigma_{k,j}$ grows, which is the mechanism driving the
acceptance of \eqref{ovo-armijo}.

\begin{assumption} \label{a4}
There is $L > 0$ such that $\|\nabla f_i(y) - \nabla f_i(x)\| \leq L \|y - x\|$
for all $i = 1, \ldots, m$ and $x, y \in \Omega$.
\end{assumption}

Assumption~\ref{a4} is the standard gradient-Lipschitz condition. It implies,
for all $i$ and $x, y \in \Omega$,
\begin{equation} \label{taylor}
f_i(y) \leq f_i(x) + \nabla f_i(x)^T(y - x) + \tfrac{L}{2}\|y - x\|^2,
\end{equation}
as well as the two-sided bound
$|f_i(y) - f_i(x) - \nabla f_i(x)^T(y - x)| \leq \tfrac{L}{2}\|y - x\|^2$. These
make $\Psi$ an upper model of $f$ once $\sigma_{k,j}$ is large relative to $L$,
which yields finite termination of the inner loop. 

\begin{lemma} \label{lem1}
Let $a_1, \ldots, a_r, b_1, \ldots, b_r \in \mathbb{R}$ and $\beta > 0$ satisfy
$a_1 \leq \cdots \leq a_r$, $b_j \leq a_j - \beta$ for $j = 1, \ldots, r$, and
$b_{i_1} \leq \cdots \leq b_{i_r}$ for a permutation $\{i_1, \ldots, i_r\}$ of
$\{1, \ldots, r\}$. Then $b_{i_q} \leq a_q - \beta$ for $q = 1, \ldots, r$.
\end{lemma}

\begin{proof}
See \cite[Lemma~2.1]{dunder1}.
\end{proof}

The next result gives a threshold on $\sigma_{k,j}$ beyond which the trial point
is accepted at Step~3, so that the inner loop terminates after finitely many
increases of $\sigma_{k,j}$, independently of $k$. The threshold scales with
$\delta^{-1}$: a narrower active band demands stronger regularization.

\begin{theorem} \label{teo:suff-descent}
Suppose that Assumptions~\ref{a1}, \ref{a3}, and \ref{a4} hold. Then, for all
$k$ and $j = 0, \ldots, j_k$, if
\begin{equation} \label{este}
\sigma_{k,j} \geq \max\{ L + 2\alpha,\ 9c_x^2/(2\delta) \},
\end{equation}
then $\xkjtrial$ satisfies \eqref{ovo-armijo}.
\end{theorem}

\begin{proof}
Assume \eqref{este}; in particular $\sigma_{k,j} \geq L + 2\alpha$. By
\eqref{taylor}, for every $i \in \{1, \ldots, m\}$,
\begin{align*}
f_i(\xkjtrial)
&\leq f_i(x^k) + \nabla f_i(x^k)^T(\xkjtrial - x^k)
+ \tfrac{L}{2}\|\xkjtrial - x^k\|^2 \\
&\leq f_i(x^k) + \nabla f_i(x^k)^T(\xkjtrial - x^k)
+ \Big( \tfrac{\sigma_{k,j}}{2} - \alpha \Big) \|\xkjtrial - x^k\|^2 \\
&= f_i(x^k) + \nabla f_i(x^k)^T(\xkjtrial - x^k)
+ \tfrac{1}{2} \big[ (\xkjtrial - x^k)^T B_{k,j} (\xkjtrial - x^k)
+ \sigma_{k,j}\|\xkjtrial - x^k\|^2 \big] \\
&\quad - \alpha \|\xkjtrial - x^k\|^2
- \tfrac{1}{2}(\xkjtrial - x^k)^T B_{k,j} (\xkjtrial - x^k).
\end{align*}
Since $B_{k,j} \succeq 0$, the last term is nonpositive and may be dropped. For
$i \in I_\delta(x^k)$, the minimizer property
$\Psi(\xkjtrial; x^k, B_{k,j}, \delta, \sigma_{k,j}) \leq
\Psi(x^k; x^k, B_{k,j}, \delta, \sigma_{k,j}) = 0$ gives
\[
\nabla f_i(x^k)^T(\xkjtrial - x^k)
+ \tfrac{1}{2} \big[ (\xkjtrial - x^k)^T B_{k,j} (\xkjtrial - x^k)
+ \sigma_{k,j}\|\xkjtrial - x^k\|^2 \big] \leq 0,
\]
so that
\begin{equation} \label{cota-lem2}
f_i(\xkjtrial) \leq f_i(x^k) - \alpha \|\xkjtrial - x^k\|^2, \qquad
i \in I_\delta(x^k).
\end{equation}
On the other hand, by \eqref{ccotasub}, the Cauchy--Schwarz inequality, and the
two-sided consequence of Assumption~\ref{a4},
$|f_i(\xkjtrial) - f_i(x^k)| \leq c_x^2/\sigma_{k,j}
+ L c_x^2/(2\sigma_{k,j}^2)$ for all $i$; using $L \leq \sigma_{k,j}$, this is
at most $3c_x^2/(2\sigma_{k,j})$, and \eqref{este} yields
\begin{equation} \label{teo7eq1}
|f_i(\xkjtrial) - f_i(x^k)| \leq \tfrac{\delta}{3}, \qquad i = 1, \ldots, m.
\end{equation}
Consequently, any index with $f_i(x^k) < f(x^k) - \delta$ satisfies
$f_i(\xkjtrial) < f(x^k)$, and any index with $f_i(x^k) > f(x^k) + \delta$
satisfies $f_i(\xkjtrial) > f(x^k)$. In particular $i_p(\xkjtrial) \in
I_\delta(x^k)$, and the number of components lying below the band is the same at
$x^k$ and $\xkjtrial$. Write $I_\delta(x^k) = \{i_1, \ldots, i_r\}$ with
$f_{i_1}(x^k) \leq \cdots \leq f_{i_r}(x^k)$, and let
$\{i'_1, \ldots, i'_r\}$ reorder them so that
$f_{i'_1}(\xkjtrial) \leq \cdots \leq f_{i'_r}(\xkjtrial)$. Since the count below
the band is preserved and $i_p(\cdot)$ lies in the band at both points,
$i_p(x^k)$ and $i_p(\xkjtrial)$ occupy the same position $q$ within the band;
that is, $i_p(x^k) = i_q$ and $i_p(\xkjtrial) = i'_q$. Applying
Lemma~\ref{lem1} with $a_j = f_{i_j}(x^k)$, $b_j = f_{i_j}(\xkjtrial)$, and
$\beta = \alpha\|\xkjtrial - x^k\|^2$ ---valid by \eqref{cota-lem2}--- gives
$f_{i'_q}(\xkjtrial) \leq f_{i_q}(x^k) - \beta$. Since
$f_{i_q}(x^k) = f_{i_p(x^k)}(x^k) = f(x^k)$,
\[
f(\xkjtrial) = f_{i_p(\xkjtrial)}(\xkjtrial) = f_{i'_q}(\xkjtrial)
\leq f(x^k) - \alpha \|\xkjtrial - x^k\|^2,
\]
which is \eqref{ovo-armijo}.
\end{proof}

The inner loop thus terminates finitely at every outer iteration. The curvature
affects only the constants below ---through the bound $M$ on $\|B_{k,j}\|$--- and
not the $O(\delta^{-2}\epsilon^{-2})$ iteration order, which coincides with that
of the first-order method \cite{alvarezovo}.

\begin{theorem} \label{teo:complexity}
Suppose that Assumptions~\ref{a1}, \ref{a3}, and \ref{a4} hold and that there is
$\flow \in \mathbb{R}$ with $f(x) \geq \flow$ for all $x \in \Omega$. Then
\begin{equation} \label{cota-sigma}
\sigma_k \leq \gamma \max\{ L + 2\alpha,\ 9c_x^2/(2\delta) \},
\end{equation}
and at most
\begin{equation} \label{cota-fcnt}
\big\lfloor 1 + \log_\gamma( \sigma_k / \sigma_{\min} ) \big\rfloor
\end{equation}
function evaluations are performed per outer iteration to obtain
\eqref{ovo-armijo}. Moreover, the number of outer iterations $k$ at which
$C_{\delta,\epsilon}$ is not satisfied by $x^{k+1}$ is at most
\begin{equation} \label{cota-iter}
\left\lfloor \frac{\big( M + \gamma \max\{ L + 2\alpha,\ 9c_x^2/(2\delta) \}
\big)^2}{\alpha} \cdot \frac{f(x^0) - \flow}{\epsilon^2} \right\rfloor.
\end{equation}
\end{theorem}

\begin{proof}
Bounds \eqref{cota-sigma} and \eqref{cota-fcnt} follow from
Theorem~\ref{teo:suff-descent} and the rule
$\sigma_{k,j+1} = \gamma\sigma_{k,j}$: starting from
$\sigma_{k,0} = \sigma_{\min}$, at most
$\lfloor 1 + \log_\gamma(\sigma_k/\sigma_{\min}) \rfloor$ increases occur before
\eqref{este} holds, after which \eqref{ovo-armijo} is satisfied.

Let $K$ be the set of outer indices $k$ at which $C_{\delta,\epsilon}$ fails at
$x^{k+1}$. Fix $k \in K$ and write $B_k = B_{k,j_k}$. By
Corollary~\ref{cor:step},
$(B_k + \sigma_k I)(x^{k+1} - x^k)
= \sum_{i=1}^n (\nu^\ell_i - \nu^u_i) e_i
- \sum_{i \in I_\delta(x^k)} \mu_i \nabla f_i(x^k)$, hence
\[
\big\| (B_k + \sigma_k I)(x^{k+1} - x^k) \big\|
= \Big\| \sum_{i \in I_\delta(x^k)} \mu_i \nabla f_i(x^k)
+ \sum_{i=1}^n (\nu^u_i - \nu^\ell_i) e_i \Big\| > \epsilon,
\]
the inequality holding because $C_{\delta,\epsilon}$ fails at $x^{k+1}$. Since
$\|B_k\| \leq M$,
\[
\epsilon < \| (B_k + \sigma_k I)(x^{k+1} - x^k) \|
\leq (M + \sigma_k) \|x^{k+1} - x^k\|,
\]
and with \eqref{cota-sigma},
$\|x^{k+1} - x^k\| > \epsilon / \big( M + \gamma\max\{L + 2\alpha, 9c_x^2/(2\delta)\} \big)$.
As $x^{k+1}$ satisfies \eqref{ovo-armijo},
\[
f(x^{k+1}) \leq f(x^k) - \alpha\|x^{k+1} - x^k\|^2
< f(x^k) - \frac{\alpha\epsilon^2}
{\big( M + \gamma\max\{L + 2\alpha, 9c_x^2/(2\delta)\} \big)^2}.
\]
Summing over $k \in K$ and using $f(x) \geq \flow$ gives
$f(x^0) - \flow > |K| \cdot \alpha\epsilon^2 /
\big( M + \gamma\max\{L + 2\alpha, 9c_x^2/(2\delta)\} \big)^2$, from which
\eqref{cota-iter} follows.
\end{proof}
%\clearpage

\begin{thebibliography}{99}

\bibitem{abmsobso} R. Andreani, E. G. Birgin, J. M. Mart\'{\i}nez and M. L. Schuverdt, Augmented Lagrangian methods under the Constant Positive Linear Dependence constraint qualification, \textit{Mathematical Programming} 111, pp. 5--32, 2008.

\bibitem{abmstango} R. Andreani, E. G. Birgin, J. M. Mart\'{\i}nez and M. L. Schuverdt, On Augmented Lagrangian methods with general lower-level constraints, \textit{SIAM Journal on Optimization} 18, pp. 1286--1309, 2008.

\bibitem{dunder1} R. Andreani, C. Dunder and J. M. Mart\'{\i}nez, Order-Value Optimization: Formulation and solution by means of a primal Cauchy method, \textit{Mathematical Methods of Operations Research} 58, pp. 387--399, 2003.

\bibitem{dunder2} R. Andreani, C. Dunder and J. M. Mart\'{\i}nez, Nonlinear-programming reformulation of the order-value optimization problem, \textit{Mathematical Methods of Operations Research} 61, pp. 365--384, 2005.

\bibitem{yano1} R. Andreani, J. M. Mart\'{\i}nez, M. Salvatierra, and F. S. Yano, Quasi-Newton methods for Order-Value Optimization and Value-at-Risk calculations, \textit{Pacific Journal of Optimization} 2, pp. 11--33, 2006.

\bibitem{yano2} R. Andreani, J. M. Mart\'{\i}nez, L. Mart\'{\i}nez, and F. S. Yano, Low order-value optimization and applications, \textit{Journal of Global Optimization} 43, pp. 1--22, 2009.

\bibitem{salva} R. Andreani, J. M. Mart\'{\i}nez, M. Salvatierra, and F. S. Yano, Global order-value optimization by means of a multistart harmonic oscillator tunneling strategy, in \textit{Global Optimization: From Theory to Implementation}, L. Liberti and N. Maculan (eds.), Springer, Boston, MA, 2006, pp. 379--404.

\bibitem{bgms} E. G. Birgin, J. L. Gardenghi, J. M. Mart\'{\i}nez, and S. A. Santos, On the use of third-order models with fourth-order regularization for unconstrained optimization, \textit{Optimization Letters} 14, pp. 815--838, 2020.

\bibitem{bgmst2} E. G. Birgin, J. L. Gardenghi, J. M. Mart\'{\i}nez, S. A. Santos, and Ph. L. Toint, Evaluation complexity for nonlinear constrained optimization using unscaled KKT conditions and high-order models, \textit{SIAM Journal on Optimization} 26, pp. 951--967, 2016.

\bibitem{bgmst1} E. G. Birgin, J. L. Gardenghi, J. M. Mart\'{\i}nez, S. A. Santos, and Ph. L. Toint, Worst-case evaluation complexity for unconstrained nonlinear optimization using high-order regularized models, \textit{Mathematical Programming} 163, pp. 359--368, 2017.

\bibitem{bmbook} E. G. Birgin and J. M. Mart\'{\i}nez, \textit{Practical Augmented Lagrangian Methods for Constrained Optimization}, Society for Industrial and Applied Mathematics, Philadelphia, 2014.

\bibitem{bmcomper} E. G. Birgin and J. M. Mart\'{\i}nez, Complexity and performance of an Augmented Lagrangian algorithm, \textit{Optimization Methods and Software} 35, pp. 885--920, 2020.

\bibitem{cogtbook} C. Cartis, N. I. M. Gould, and Ph. L. Toint, \textit{Evaluation Complexity of Algorithms for Nonconvex Optimization: Theory, Computation, and Perspectives}, SIAM, Philadelphia, PA, 2022.

\bibitem{farrington} C. P. Farrington, Modelling forces of infection for measles, mumps and rubella, \textit{Statistics in Medicine} 9, pp. 953--967, 1990.

\bibitem{gauvin} J. Gauvin, A necessary and sufficient regularity condition to have bounded multipliers in nonconvex programming, \textit{Mathematical Programming} 12, pp. 136--138.
  
\bibitem{jorion} P. Jorion, \textit{Value at {R}isk: {T}he new benchmark for managing financial risk}, 3rd. ed., McGraw-Hill, 2009.

\bibitem{mtop} J. M. Mart\'{\i}nez, Generalized Order-Value Optimization, \textit{TOP} 20, pp. 75--98, 2012.

\bibitem{mgh} J. J. Mor\'e, B. S. Garbow, and K. E. Hillstrom, Testing unconstrained optimization software, \textit{ACM Transactions on Mathematical Software} 7, pp. 17--41, 1981.
  
\bibitem{nesterov} Y. Nesterov and B. T. Polyak, Cubic regularization of Newton method and its global performance, \textit{Mathematical Programming} 108, 177--205, 2006.

\bibitem{jiang} Z. Jiang, Q. Hu, and X. Zheng, Optimality condition and complexity of order-value optimization problems and low order-value optimization problems, \textit{Journal of Global Optimization} 69, pp. 511--523, 2017.

\bibitem{alvarezovo} G. Q. Álvarez and E. G. Birgin, A first-order regularized approach to the order-value optimization problem, \textit{Optimization Methods and Software} 40, pp. 650--674, 2025.
    
\end{thebibliography}

\end{document}
